%% ============================================================
%% BÖLÜM 16 — Varyans Analizi (ANOVA)
%% Olasılık Teorisi ve Matematiksel İstatistik
%% ============================================================
\chapter{Varyans Analizi (ANOVA)}
\label{chp:anova}

\bolumalinti{%
  The analysis of variance is not a mathematical theorem, but rather a convenient method of arranging the arithmetic.
}{Ronald A. Fisher}

\begin{kazanimlar}
Bu bölümün sonunda okuyucu:
\begin{itemize}
  \item Tek yönlü ANOVA modelini kurar; $F$-testini türetir ve uygular.
  \item Toplam kareler toplamı ayrışımını $\mathrm{SST} = \mathrm{SSA} + \mathrm{SSE}$ ispat eder.
  \item İki yönlü ANOVA'da ana etkiler ve etkileşim kavramlarını açıklar.
  \item Çoklu karşılaştırma yöntemlerini (Tukey, Bonferroni) uygular.
  \item Sabit, rasgele ve karma etkili modelleri ayırt eder.
  \item ANOVA ile çok değişkenli regresyon arasındaki ilişkiyi gösterir.
\end{itemize}
\end{kazanimlar}

\begin{motivasyon}
İki grubun ortalamasını karşılaştırmak için $t$-testi vardır. Üç veya daha fazla gruba genişletebilir miyiz? Tek tek $t$-testleri uygulamak tip I hata oranını patlatır. ANOVA gruplar arasındaki varyasyon ile grup içi varyasyonu karşılaştırarak bu sorunu çözer.

\medskip
\noindent\textbf{Köprü:} Bölüm~\ref{chp:dogrusal_modeller}'deki regresyon modeli ve $F$-testi, ANOVA'nın özel bir durumudur — $\mathbf{X}$ tasarım matrisi gösterge (dummy) değişkenlerden oluştuğunda ANOVA elde edilir.
\end{motivasyon}

%% ============================================================
\section{Tek Yönlü ANOVA}
\label{sec:tek_anova}
%% ============================================================

\begin{tanim}[Tek yönlü ANOVA modeli]\label{def:tek_anova}
$k$ grup; $i$. grupta $n_i$ gözlem:
\[
  Y_{ij} = \mu + \alpha_i + \varepsilon_{ij}, \quad i=1,\ldots,k;\; j=1,\ldots,n_i,
\]
$\varepsilon_{ij}\iid\Normal(0,\sigma^2)$, $\sum n_i\alpha_i=0$ (kısıt). $N=\sum n_i$ toplam gözlem. Hipotez:
\[
  H_0: \alpha_1=\alpha_2=\cdots=\alpha_k=0 \quad\text{(tüm grup ortalamaları eşit)}.
\]
\end{tanim}

\subsection{Kareler Toplamı Ayrışımı}

\begin{teorem}[SST ayrışımı]\label{thm:sst_ayrisim}
\[
  \underbrace{\sum_{i=1}^k\sum_{j=1}^{n_i}(Y_{ij}-\bar{Y}_{..})^2}_{\mathrm{SST}} =
  \underbrace{\sum_{i=1}^k n_i(\bar{Y}_{i.}-\bar{Y}_{..})^2}_{\mathrm{SSA}} +
  \underbrace{\sum_{i=1}^k\sum_{j=1}^{n_i}(Y_{ij}-\bar{Y}_{i.})^2}_{\mathrm{SSE}}.
\]
\end{teorem}

\begin{proof}
$Y_{ij}-\bar Y_{..} = (\bar Y_{i.}-\bar Y_{..}) + (Y_{ij}-\bar Y_{i.})$ ayrışımını kareye al ve çift çarpım teriminin $\sum_{i,j}(\bar Y_{i.}-\bar Y_{..})(Y_{ij}-\bar Y_{i.})=0$ olduğunu gör (her $i$ için $\sum_j(Y_{ij}-\bar Y_{i.})=0$).
\end{proof}

\begin{teorem}[ANOVA $F$-testi]\label{thm:anova_f}
$H_0$ altında:
\[
  F = \frac{\mathrm{SSA}/(k-1)}{\mathrm{SSE}/(N-k)} \sim F_{k-1,N-k}.
\]
$H_0$ reddedilir ise $F > F_{k-1,N-k,\alpha}$.
\end{teorem}

\begin{proof}
$\mathrm{SSE}/\sigma^2\sim\chi^2_{N-k}$ (her grup içi $\chi^2_{n_i-1}$; bağımsız toplam). $H_0$ altında $\mathrm{SSA}/\sigma^2\sim\chi^2_{k-1}$ (ortogonal projeksiyon). SSA $\perp$ SSE. $F$ tanımı gereği $F_{k-1,N-k}$.
\end{proof}

\begin{ornek}[ANOVA tablosu]\label{ex:anova_tablo}
\[
\begin{array}{lcccc}
\hline
\text{Kaynak} & \text{SD} & \text{KT} & \text{MKT} & F \\
\hline
\text{Gruplar} & k-1 & \mathrm{SSA} & \mathrm{MSA}=\mathrm{SSA}/(k-1) & \mathrm{MSA}/\mathrm{MSE} \\
\text{Hata} & N-k & \mathrm{SSE} & \mathrm{MSE}=\mathrm{SSE}/(N-k) & \\
\text{Toplam} & N-1 & \mathrm{SST} & & \\
\hline
\end{array}
\]
SD: serbestlik derecesi; KT: kareler toplamı; MKT: ortalama KT.
\end{ornek}

\begin{teorem}[Beklenti MSE ve MSA]\label{thm:mse_msa_beklenti}
\[
  \Beklenti[\mathrm{MSE}] = \sigma^2, \qquad
  \Beklenti[\mathrm{MSA}] = \sigma^2 + \frac{\sum n_i\alpha_i^2}{k-1}.
\]
\end{teorem}

\begin{proof}
MSE yansız ($\Beklenti[\mathrm{SSE}]=(N-k)\sigma^2$). MSA: $\Beklenti[\mathrm{SSA}]=\sum n_i\Var(\bar Y_{i.}-\bar Y_{..})+(k-1)\sigma^2+\sum n_i\alpha_i^2 = (k-1)\sigma^2+\sum n_i\alpha_i^2$.
\end{proof}

%% ============================================================
\section{Çoklu Karşılaştırma}
\label{sec:coklu_karsilastirma}
%% ============================================================

\begin{kesif}
$F$-testi $H_0$'ı reddedince hangi grupların farklı olduğunu söylemez. Post-hoc testler bu soruyu yanıtlar ama FWER kontrolü gereklidir.
\end{kesif}

\begin{tanim}[Tukey HSD]\label{def:tukey}
Dengeli tasarımda ($n_i=n$ herkese eşit). \textbf{Tukey dürüst anlamlı fark (HSD)}:
\[
  |\bar{Y}_{i.}-\bar{Y}_{j.}| > q_{\alpha,k,N-k}\!\sqrt{\frac{\mathrm{MSE}}{n}},
\]
$q_{\alpha,k,N-k}$ studentleştirilmiş aralık ($q$) dağılımının kritik değeri. Tüm $\binom{k}{2}$ çift karşılaştırmada FWER $=\alpha$.
\end{tanim}

\begin{teorem}[Scheffé yöntemi]\label{thm:scheffe}
Herhangi kontrast $L=\sum c_i\mu_i$ ($\sum c_i=0$) için $\%100(1-\alpha)$ güven aralığı:
\[
  \sum c_i\bar{Y}_{i.} \pm \sqrt{(k-1)F_{k-1,N-k,\alpha}\cdot\mathrm{MSE}\sum\frac{c_i^2}{n_i}}.
\]
Tüm olası kontrastları kapsayan en geniş yöntem; bireysel karşılaştırmalarda muhafazakâr.
\end{teorem}

%% ============================================================
\section{İki Yönlü ANOVA}
\label{sec:iki_anova}
%% ============================================================

\begin{tanim}[İki yönlü ANOVA modeli]\label{def:iki_anova}
$a$ düzeyli Faktör A, $b$ düzeyli Faktör B; her hücrede $n$ gözlem:
\[
  Y_{ijk} = \mu + \alpha_i + \beta_j + (\alpha\beta)_{ij} + \varepsilon_{ijk},
\]
$i=1,\ldots,a$; $j=1,\ldots,b$; $k=1,\ldots,n$; $\varepsilon_{ijk}\iid\Normal(0,\sigma^2)$.

Kısıtlar: $\sum\alpha_i=\sum\beta_j=\sum_i(\alpha\beta)_{ij}=\sum_j(\alpha\beta)_{ij}=0$.
\end{tanim}

\begin{teorem}[İki yönlü SST ayrışımı]\label{thm:iki_sst}
\[
  \mathrm{SST} = \mathrm{SSA} + \mathrm{SSB} + \mathrm{SSAB} + \mathrm{SSE}.
\]
Her terim karşılık gelen serbestlik derecesiyle $F$-testi oluşturur.
\end{teorem}

\begin{ornek}[İki yönlü ANOVA tablosu]\label{ex:iki_anova_tablo}
\[
\begin{array}{lccc}
\hline
\text{Kaynak} & \text{SD} & F \\
\hline
A & a-1 & \mathrm{MSA}/\mathrm{MSE} \\
B & b-1 & \mathrm{MSB}/\mathrm{MSE} \\
A\times B & (a-1)(b-1) & \mathrm{MSAB}/\mathrm{MSE} \\
Hata & ab(n-1) & \\
\hline
\end{array}
\]
\end{ornek}

\begin{tanim}[Etkileşim]\label{def:etkilesim}
$(\alpha\beta)_{ij}\neq0$: \textbf{etkileşim} \emph{(interaction)} mevcut. Faktör A'nın etkisi Faktör B'nin düzeyine göre değişiyor demektir. Etkileşim varken ana etkilerin yorumu anlamsız olabilir.
\end{tanim}

\begin{kritiksorgu}
Etkileşim anlamlıyken ana etkilere bakılmalı mı? Genellikle etkileşim istatistiği önce kontrol edilir; anlamlıysa ana etkiler ayrı ayrı her faktörün düzeyinde incelenir.
\end{kritiksorgu}

%% ============================================================
\section{Rasgele Etkili Model}
\label{sec:rasgele_etkili}
%% ============================================================

\begin{tanim}[Tek yönlü rasgele etki modeli]\label{def:rasgele_etki}
$\alpha_i\iid\Normal(0,\sigma_A^2)$ (sabit değil rasgele). Model:
\[
  Y_{ij} = \mu + \alpha_i + \varepsilon_{ij},
\]
$\alpha_i\perp\varepsilon_{ij}$. $\Var(Y_{ij})=\sigma_A^2+\sigma^2$; $\Cov(Y_{ij},Y_{ij'})=\sigma_A^2$ ($j\neq j'$).

\textbf{İçsınıf korelasyon katsayısı:} $\mathrm{ICC} = \sigma_A^2/(\sigma_A^2+\sigma^2)$.
\end{tanim}

\begin{teorem}[Rasgele etki modeli $F$-testi]\label{thm:rasgele_f}
$H_0:\sigma_A^2=0$ altında:
\[
  F = \frac{\mathrm{MSA}}{\mathrm{MSE}} \sim F_{k-1,N-k}.
\]
$F$-istatistiği aynı; yorumu farklı: gruplar arası varyans sıfır mı?
\end{teorem}

\begin{ispatiskele}
$H_0:\sigma_A^2=0$ altında $\Beklenti[\mathrm{MSA}]=\Beklenti[\mathrm{MSE}]=\sigma^2$; yani $F$ ortalaması $\approx 1$. $H_0$ doğruyken $\mathrm{SSA}/\sigma^2\sim\chi^2_{k-1}$ (rasgele $\alpha_i$ katkısı yoktur). Aynı $F$ dağılımı.
\end{ispatiskele}

%% ============================================================
\section{ANOVA — Regresyon Bağlantısı}
\label{sec:anova_regresyon}
%% ============================================================

\begin{teorem}[ANOVA = regresyon]\label{thm:anova_regresyon}
$k$ grup tek yönlü ANOVA, $k-1$ gösterge değişkenli regresyona eşdeğerdir:
\[
  x_{i\ell} = \mathbf{1}[i=\ell], \quad \ell=1,\ldots,k-1.
\]
ANOVA $F$-istatistiği ile regresyon $F$-istatistiği özdeştir.
\end{teorem}

\begin{proof}
Tasarım matrisi $\mathbf{X}$ (intercept + $k-1$ gösterge sütunu): $\hat{\boldsymbol\beta}=(\mu,\alpha_1-\alpha_k,\ldots,\alpha_{k-1}-\alpha_k)^\top$. $\mathrm{SSR}=\mathrm{SSA}$, $\mathrm{SSE}$ özdeş. $F$ aynı.
\end{proof}

\begin{disiplinlerarasibaglan}
\textbf{Tıp:} Klinik deneylerde tedavi grupları karşılaştırması. \textbf{Eğitim araştırmaları:} Farklı öğretim yöntemlerinin etkinliği. \textbf{Tarımsal araştırma:} Fisher'in orijinal motivasyonu — farklı gübre miktarlarının mahsule etkisi.
\end{disiplinlerarasibaglan}

%% ============================================================

\begin{mikroozet}
\begin{itemize}
  \item Tek yönlü ANOVA: SST = SSA + SSE; $F = \mathrm{MSA}/\mathrm{MSE}\sim F_{k-1,N-k}$.
  \item Tukey HSD: dengeli tasarımda tüm çift karşılaştırmalarda FWER kontrolü.
  \item İki yönlü ANOVA: etkileşim terimi ana etkilerin yorumunu etkiler; önce kontrol et.
  \item Rasgele etki: $\sigma_A^2$ varyans bileşeni; ICC; aynı $F$ farklı yorum.
\end{itemize}
\end{mikroozet}

\begin{ozyansima}
ANOVA normallik ve eşit varyans gerektirir. Bu varsayımlar bozulduğunda ne yapılır? Kruskal-Wallis testi bir alternatiftir ama gücü ne kadar?
\end{ozyansima}

\begin{kopru}
\textbf{Nereden geldik:} Bölüm~\ref{chp:dogrusal_modeller}'deki regresyon ve $F$-testi; Bölüm~\ref{chp:hipotez_testi}'deki çoklu test düzeltmesi.

\textbf{Bu kitabın sonu:} Olasılık teorisinden istatistiksel çıkarıma uzanan yolculuğu tamamladık. Bundan sonraki adımlar: Bölüm~\ref{chp:stokastik} ve~\ref{chp:asimptotik}'daki ileri konular, semiparametrik model ve Bayes hesaplamalı yöntemler.
\end{kopru}

%% ============================================================
%% ALISTIRMALAR
%% ============================================================

\cozumlualistirmalar

\begin{alistirma}[Al.16.1]{A}{Tek yönlü ANOVA hesabı}
3 grup: A$=\{4,6,8\}$, B$=\{3,5,7\}$, C$=\{8,10,12\}$. ANOVA tablosunu oluştur ve $\alpha=0.05$'te $H_0$ test et.
\end{alistirma}
\alcozum{%
$\bar Y_{A.}=6$, $\bar Y_{B.}=5$, $\bar Y_{C.}=10$; $\bar Y_{..}=7$. SSA$=3[(6-7)^2+(5-7)^2+(10-7)^2]=3[1+4+9]=42$. SSE$=[(4-6)^2+(6-6)^2+(8-6)^2]+[(3-5)^2+(5-5)^2+(7-5)^2]+[(8-10)^2+(10-10)^2+(12-10)^2]=8+8+8=24$. $F=(42/2)/(24/6)=21/4=5.25$. $F_{2,6,0.05}=5.14$. $F>5.14$: $H_0$ reddedilir.%
}

\begin{alistirma}[Al.16.2]{B}{SST ayrışımı ispatı}
Genel formel: $(Y_{ij}-\bar Y_{..}) = (\bar Y_{i.}-\bar Y_{..}) + (Y_{ij}-\bar Y_{i.})$. Her iki tarafın karesini alıp $i,j$ üzerinden toplayınca çift çarpım teriminin kaybolduğunu göster.
\end{alistirma}
\alcozum{%
$\sum_{i,j}(\bar Y_{i.}-\bar Y_{..})(Y_{ij}-\bar Y_{i.}) = \sum_i(\bar Y_{i.}-\bar Y_{..})\sum_j(Y_{ij}-\bar Y_{i.})$. Her $i$ için iç toplam $\sum_j(Y_{ij}-\bar Y_{i.})=0$. Dolayısıyla çift çarpım $=0$; SST = SSA + SSE.%
}

\alistirmalar

\begin{alistirma}[Al.16.3]{A}{Tukey HSD}
Alıştırma~\ref{Al.16.1}'deki veriler için Tukey HSD yöntemiyle (\%95) hangi grup çiftlerinin anlamlı biçimde farklı olduğunu belirle. ($q_{0.05,3,6}\approx4.34$.)
\end{alistirma}
\alcozum{%
MSE$=4$; $n=3$. HSD $=4.34\sqrt{4/3}=4.34\cdot1.155=5.01$. $|A-B|=1<5.01$, $|A-C|=4<5.01$, $|B-C|=5<5.01$ — Hiç çift anlamlı değil (görece küçük örneklem). Not: $F$-testi anlamlı bulmuştu ama Tukey daha muhafazakâr.%
}

\begin{alistirma}[Al.16.4]{B}{İki yönlü ANOVA — etkileşim}
$a=2$ (A: erkek/kadın), $b=2$ (B: ilaç/plasebo), $n=4$. Hücre ortalamaları: $\bar Y_{11}=20$, $\bar Y_{12}=15$, $\bar Y_{21}=22$, $\bar Y_{22}=25$. MSE$=10$, $N=16$. Etkileşim F-istatistiğini hesapla.
\end{alistirma}
\alcozum{%
$\bar Y_{1.}=17.5$, $\bar Y_{2.}=23.5$, $\bar Y_{.1}=21$, $\bar Y_{.2}=20$; $\bar Y_{..}=20.5$. SSAB $= 4\sum_{ij}[(Y_{ij}-\bar Y_{i.}-\bar Y_{.j}+\bar Y_{..})^2]$. Hücre etkileşim: $(\alpha\beta)_{11}=20-17.5-21+20.5=2$; $(\alpha\beta)_{12}=15-17.5-20+20.5=-2$; $(\alpha\beta)_{21}=22-23.5-21+20.5=-2$; $(\alpha\beta)_{22}=25-23.5-20+20.5=2$. SSAB$=4(4+4+4+4)=64$. MSAB$=64/1=64$. $F=64/10=6.4$. $F_{1,12,0.05}\approx4.75$: etkileşim anlamlı.%
}

\begin{alistirma}[Al.16.5]{C}{Rasgele etki: ICC tahmin}
$k=5$ sınıf, her sınıftan $n=10$ öğrenci test skoru. MSA$=120$, MSE$=30$. $\hat\sigma_A^2$ ve $\hat{\mathrm{ICC}}$'yi hesapla. Yorumla.
\end{alistirma}
\alcozum{%
$\Beklenti[\mathrm{MSA}]=\sigma^2+n\sigma_A^2$; $\hat\sigma^2=\mathrm{MSE}=30$; $\hat\sigma_A^2=(\mathrm{MSA}-\mathrm{MSE})/n=(120-30)/10=9$. $\hat{\mathrm{ICC}}=9/(9+30)=9/39\approx0.23$. Yorum: sınıf üyeliği test skoru varyansının \%23'ünü açıklıyor; sınıf etkisi önemli.%
}

\begin{alistirma}[Al.16.6]{C}{ANOVA — regresyon eşdeğerliği}
$k=3$ grup ANOVA'yı gösterge değişkenli regresyon olarak yeniden kur. Tasarım matrisini yaz; EKK tahminlerinin $\hat\mu$, $\hat\alpha_1-\hat\alpha_3$, $\hat\alpha_2-\hat\alpha_3$ olduğunu göster.
\end{alistirma}
\alcozum{%
$\mathbf{X}=[\mathbf{1},\mathbf{x}_1,\mathbf{x}_2]$; $x_{i1}=\mathbf{1}[$grup 1$]$, $x_{i2}=\mathbf{1}[$grup 2$]$. Normal denklem çözümü: intercept $\hat\beta_0=\bar Y_{3.}=\hat\mu+\hat\alpha_3$; $\hat\beta_1=\bar Y_{1.}-\bar Y_{3.}=\hat\alpha_1-\hat\alpha_3$; $\hat\beta_2=\bar Y_{2.}-\bar Y_{3.}=\hat\alpha_2-\hat\alpha_3$. $F=\mathrm{SSR}/2\div\mathrm{SSE}/(N-3)$ = ANOVA $F$.%
}
